\documentclass[12pt]{book}
\usepackage[utf8]{inputenc}
\usepackage[T1]{fontenc}
\usepackage{amsmath,amssymb,amsfonts}
\usepackage{amsthm}

% Standard operators
\DeclareMathOperator{\Spec}{Spec}
\DeclareMathOperator{\Hom}{Hom}
\DeclareMathOperator{\Res}{Res}

\begin{document}

\setcounter{page}{363}
that \(a_{i}=0\) for \(i<-r\). Then \(Z = \Spec \sheaf{O}_{x}/t^{r}\) will do.
Our element \(u\) is then identified with the element
\[\in \Hom_{\sheaf{O}_{x}}\left(B, I \otimes_{A} \omega_{x} \otimes K / \sheaf{O}_{x}\right)\]
defined by \(\psi(1)=u\). This corresponds by the residue isomorphism to a certain
\[\varphi \in \Hom_{A}(B,I),\]
and by definition of the trace morphism we then have
\[\Gamma\left(Tr_{f,x}\right)(u)=\varphi(1) .\]

But by the lemma above, \(\varphi(1)\) is the coefficient of \(t^{-1} dt\), so our trace is \(a_{-1}\), as required.

\setcounter{page}{364}
Remark. Since the trace morphism has already been defined, and we are here calculating it, we can state as a corollary that the process of taking the coefficient of \(t^{-1} dt\) in \(u\) is independent of the local parameter \(t\) chosen.
This is in strong contrast with the approach of Serre's book [16, Ch. II] where the residue of a differential is defined as \(a_{-1}\), and a tedious proof is required to show that it is independent of the local parameter [loc. cit. Prop. 5].

To establish the relation between his approach and ours, we make the following

Definition. Let \(M\) be an \(A\)-module, and let
\[w \in M \otimes_{A} \omega_{\eta}\]
be a meromorphic differential on \(X\).
Then for each closed point \(x \in X\), rational over \(Y\), we define the residue of \(w\) at \(x\), as follows.
Choose a local parameter \(t\) at \(x\). Let
\[w_{x} \in M \otimes_{A} \omega_{x} \otimes K / \sheaf{O}_{x}\]
be the image of \(w\) under the natural map
\[\omega_{\eta}=\omega_{x} \otimes K \to \omega_{x} \otimes K / \sheaf{O}_{x} .\]

Write
\[w_{x}=\sum_{i<0} a_{i} t^{i} dt\]
with \(a_{i} \in M\), and then define
\[\Res_{x}(w)=a_{-1} .\]

\setcounter{page}{365}
(Observe that \(a_{-1}\) is independent of the choice of local parameter. Indeed, we may assume that \(M\) is finitely generated, in which case it is a submodule of \(I^{n}\) for some \(n\), so we reduce to the case \(M=I\), where \(a_{-1}=\Gamma(Tr_{f,x})\) which is independent of \(t\).)

Lemma. Residues have the following properties
\begin{enumerate}
\item \(\Res_{x}(w)\) is \(A\)-linear in \(w\).
\item \(\Res_{x}(w)=0\) if \(w \in M \otimes \omega_{x}\).
\item \(\Res_{x}(dg)=0\) for any \(g \in M \otimes_{A} K\).
\item \(\Res_{x}(t^{-1} dt)=1\) if \(t\) is a local parameter at \(x\).
\end{enumerate}

Proof. All properties follow from the definition.

\setcounter{page}{366}
Theorem 1.5
Let \(A\) be a local Artin ring with algebraically closed residue field \(k\), let \(Y=\Spec A\), and let \(X=\mathbb{P}_{Y}^{1}\) be the projective line over \(Y\).
Let \(M\) be an \(A\)-module, and let
\[w \in M \otimes_{A} \omega_{\eta}\]
be a meromorphic differential on \(X\).
Then
\[\sum_{\substack{x \in X \\ closed }} \Res_{x}(w)=0\]

Proof. (Modeled on the proof of [16, Ch. II, Lemma 3].)
First of all, we may assume that \(M\) is finitely generated.
Then writing \(M\) as a quotient of \(A^{n}\) for suitable \(n\), we reduce to the case \(M=A\).
Let \(w \in \omega_{\eta}\) be a meromorphic differential.
Let \(X_0\) be the affine line contained as an open subset of \(X\).
Then from the exact sequence of sheaves
\[0 \to \omega \to i_{\eta}(\omega_{\eta}) \to \bigoplus_{x \in X} i_{x}\left(\omega_{x} \otimes K / \sheaf{O}_{x}\right) \to 0\]
we obtain the following exact sequence of global sections on \(X_0\):
\[0 \to \Gamma(X_0,\omega) \to \omega_{\eta} \to \bigoplus_{x \in X_0} \omega_{x} \otimes K / \sheaf{O}_{x} \to 0 .\]

\setcounter{page}{367}
For each closed point \(x \in X_0\), we can choose a local parameter of the form \(t-c_{x}\), with \(c_{x} \in A\).
Then our exact sequence shows that \(w\) can be written uniquely in the form
\[w=f(t) dt+\sum_{\substack{x \in X_0 \\ closed }} \sum_{i<0} a_{i,x}\left(t-c_{x}\right)^{i} dt\]
where \(f(t) \in A[t]\) is a polynomial in \(t\) with coefficients in \(A\).
By linearity, we reduce to proving the residue formula in the two cases
\[(1)\quad w=t^{n} dt,\quad n \ge 0\]
\[(2)\quad w=\left(t-c_{x_0}\right)^{n},\quad n<0.\]

In the first case, \(\Res_{x}(w)=0\) for every \(x \in X_0\).
We take \(s=1/t\) as a local point at \(\infty\).
Then \(w=-s^{-n-2} ds\), so the residue is zero.

In the second case, for any \(x \ne x_0\), \(t-c_{x_0}\) is invertible, so its residue is zero.
Take \(s=1/(t-c_{x_0})\) as local parameter at \(\infty\).
Then \(w=-s^{-n-2} ds\), residue zero.
Hence the sum of the residues is zero. q.e.d.

\setcounter{page}{368}
Remark. It follows from the theorem of the next section that the theorem remains true for any smooth proper curve \(X\) over \(Y\).

Corollary 1.6.
\[Tr_{f}: f_{*} f^{\Delta}(\tilde{I}) \to \tilde{I}\]
is a morphism of complexes.

Proof. Using Proposition 1.3 above, this is precisely the statement of the Theorem in the case \(M = I\).

\end{document}